%%%%%%%%%%%%%%%%%%%%%%%%%%%%%%%%%%%%%%%%%
% Short Sectioned Assignment LaTeX Template Version 1.0 (5/5/12)
% This template has been downloaded from: http://www.LaTeXTemplates.com
% Original author:  Frits Wenneker (http://www.howtotex.com)
% License: CC BY-NC-SA 3.0 (http://creativecommons.org/licenses/by-nc-sa/3.0/)
%%%%%%%%%%%%%%%%%%%%%%%%%%%%%%%%%%%%%%%%%

% \documentclass[paper=a4, fontsize=11pt]{scrartcl} % A4 paper and 11pt font size
\documentclass[11pt, a4paper]{book}
\usepackage[T1]{fontenc} % Use 8-bit encoding that has 256 glyphs
\usepackage[utf8]{inputenc}
\usepackage{fourier} % Use the Adobe Utopia font for the document - comment this line to return to the LaTeX default
\usepackage{listings} % para insertar código con formato similar al editor
\usepackage[spanish, es-tabla]{babel} % Selecciona el español para palabras introducidas automáticamente, p.ej. "septiembre" en la fecha y especifica que se use la palabra Tabla en vez de Cuadro
\usepackage{url} % ,href} %para incluir URLs e hipervínculos dentro del texto (aunque hay que instalar href)
\usepackage{graphics,graphicx, float} %para incluir imágenes y colocarlas
\usepackage[gen]{eurosym} %para incluir el símbolo del euro
\usepackage{cite} %para incluir citas del archivo <nombre>.bib
\usepackage{enumerate}
\usepackage{hyperref}
\usepackage{graphicx}
\usepackage{tabularx}
\usepackage{booktabs}

\usepackage[table,xcdraw]{xcolor}
\hypersetup{
	colorlinks=true,	% false: boxed links; true: colored links
	linkcolor=black,	% color of internal links
	urlcolor=cyan		% color of external links
}
\renewcommand{\familydefault}{\sfdefault}
\usepackage{fancyhdr} % Custom headers and footers
\pagestyle{fancyplain} % Makes all pages in the document conform to the custom headers and footers
\fancyhead[L]{} % Empty left header
\fancyhead[C]{} % Empty center header
\fancyhead[R]{Me Llamo Así} % My name
\fancyfoot[L]{} % Empty left footer
\fancyfoot[C]{} % Empty center footer
\fancyfoot[R]{\thepage} % Page numbering for right footer
%\renewcommand{\headrulewidth}{0pt} % Remove header underlines
\renewcommand{\footrulewidth}{0pt} % Remove footer underlines
\setlength{\headheight}{13.6pt} % Customize the height of the header

\usepackage{titlesec, blindtext, color}
\definecolor{gray75}{gray}{0.75}
\newcommand{\hsp}{\hspace{20pt}}
\titleformat{\chapter}[hang]{\Huge\bfseries}{\thechapter\hsp\textcolor{gray75}{|}\hsp}{0pt}{\Huge\bfseries}
\setcounter{secnumdepth}{4}
\usepackage[Lenny]{fncychap}


\begin{document}

	% Plantilla portada UGR
	\begin{titlepage}
\newlength{\centeroffset}
\setlength{\centeroffset}{-0.5\oddsidemargin}
\addtolength{\centeroffset}{0.5\evensidemargin}
\thispagestyle{empty}

\noindent\hspace*{\centeroffset}\begin{minipage}{\textwidth}

\centering
\includegraphics[width=0.9\textwidth]{logos/logo_ugr.jpg}\\[1.4cm]

\textsc{ \Large TRABAJO FIN DE GRADO\\[0.2 cm]}
\textsc{ GRADO EN INGENIERÍA INFORMÁTICA}\\[1 cm]

{\Huge\bfseries Título \\}
\noindent\rule[-1ex]{\textwidth}{3pt}\\[3.5 ex]
{\large\bfseries Subtítulo }
\end{minipage}

\vspace{2.5cm}
\noindent\hspace*{\centeroffset}
\begin{minipage}{\textwidth}
\centering

\textbf{Autor}\\ {Estudiante}\\[2.5 ex]
\textbf{Director}\\ {Tutor(a)(es)}\\[2 cm]
\includegraphics[width=0.3\textwidth]{logos/etsiit_logo.png}\\[0.1 cm]
\textsc{Escuela Técnica Superior de Ingenierías Informática y de Telecomunicación}\\
\textsc{---}\\
Granada, junio de 2023
\end{minipage}
\end{titlepage}


	% Plantilla prefacio UGR
	\thispagestyle{empty}

\begin{center}
{\large\bfseries Título \\ Subtítulo }\\
\end{center}
\begin{center}
Nombre Del Estudiante\\
\end{center}

%\vspace{0.7cm}

\vspace{0.5cm}
\noindent\textbf{Palabras clave}: \textit{software libre}
\vspace{0.7cm}

\noindent\textbf{Resumen}\\
	

\cleardoublepage

\begin{center}
	{\large\bfseries Same, but in English}\\
\end{center}
\begin{center}
	Student's name\\
\end{center}
\vspace{0.5cm}
\noindent\textbf{Keywords}: \textit{open source}, \textit{floss}
\vspace{0.7cm}

\noindent\textbf{Abstract}\\


\cleardoublepage

\thispagestyle{empty}

\noindent\rule[-1ex]{\textwidth}{2pt}\\[4.5ex]

D. \textbf{Tutora/e(s)}, Profesor(a) del ...

\vspace{0.5cm}

\textbf{Informo:}

\vspace{0.5cm}

Que el presente trabajo, titulado \textit{\textbf{Chief}},
ha sido realizado bajo mi supervisión por \textbf{Estudiante}, y autorizo la defensa de dicho trabajo ante el tribunal
que corresponda.

\vspace{0.5cm}

Y para que conste, expiden y firman el presente informe en Granada a Junio de 2018.

\vspace{1cm}

\textbf{El/la director(a)/es: }

\vspace{5cm}

\noindent \textbf{(nombre completo tutor/a/es)}

\chapter*{Agradecimientos}






	% Índice de contenidos
	\newpage
	\tableofcontents

	% Índice de imágenes y tablas
	\newpage
	\listoffigures

	% Si hay suficientes se incluirá dicho índice
	\listoftables 
	\newpage

	% Introducción 
	\chapter{Introducción}

Un tropo es una figura retórica o convención que se usa en la creación de historias y es reconocible por la audiencia. Un tropo puede ser una situación, un tipo de personaje, una estructura narrativa o cualquier patrón que se utilice a la hora de contar una historia, de forma que no la interrumpa (en cuyo caso, sería un cliché) sino que ayuda a entenderla y construirla mejor. Los tropos tienen interés como una parte muy importante del puzle que es producir una obra de cualquier tipo, ya sea una serie, película, videojuego, libro... Estos se pueden analizar para obtener una serie de pautas que puedan servir para crear nuevas obras más efectivas, ya sea analizando cómo evoluciona el uso de tropos a lo largo del tiempo y cuáles son los más adecuados para representar un tipo de género o historia concreta. En resumen, los tropos son clave en la popularidad que puede llegar a tener una obra. 

TvTropes \cite{tvtropes} es una wiki que describe y almacena ejemplos de tropos en cualquier tipo de medio audiovisual. Contiene una gran lista de tropos, con una descripción y ejemplos reales de dónde y cómo se aplican. TvTropes indica para cada entrada de una obra audiovisual todos los tropos que se le han identificado y por qué, con ejemplos, siendo esta la principal fuente de información de la web. Esta información está constantemente creciendo.

Sin embargo, tal y como se define en la propia web, el objetivo de TvTropes es ser una wiki sobre tropos que se usan para contar historias y mostrarlos de un modo accesible y divertido de leer, destinado al usuario final. Esta información no está estructurada y, por tanto, no existe una manera de acceder a ella (mediante una API, por ejemplo) para su uso en ciencia de datos o inteligencia artificial. Además de esto, TvTropes es una web centrada en los tropos y nada más que los tropos, por lo que no manejan información secundaria sobre los contenidos audiovisuales que tienen, los cuales pueden ser también interesantes de analizar.

En estos análisis entran trabajos de investigación como el de \cite{garcia2020tropes} o el de \cite{tfgantoniogamiz}, los cuales actúan como principal motivación de este proyecto, puesto que ambos utilizan tanto una base de datos existente que está desactualizada como una biblioteca de scraping de TvTropes que no consigue extraer toda la información que tiene la web actualmente, la cual ha cambiado bastante con el tiempo.

Por tanto, ante la necesidad de tener toda la información completa de TvTropes, se quiere construir un scraper capaz de obtener de forma estructurada esta información sobre tropos, ampliando aún más su funcionalidad para obtener tropos de cualquier medio audiovisual e información de metadatos adicional, para que cualquier programador, científico de datos o investigador pueda emplearla para crear modelos, programas o investigaciones.\\

Este trabajo se divide en tres partes diferenciadas. La primera empezará describiendo el problema que se desea resolver y los objetivos que se quieren alcanzar, para luego analizar las alternativas que existen actualmente en materia de scrapers y herramientas que extraen la información de la web de TvTropes, así como la situación actual de la propia página y los retos que presenta a la hora de extraer su información. La segunda parte se centrará en explicar y desarrollar la metodología de desarrollo ágil que se ha utilizado y cómo se ha guiado todo el desarrollo del software hasta el final. Por último, la tercera parte desarrollará todo el trabajo realizado y cómo se ha llegado a la solución final, valorando finalmente si se han alcanzado los objetivos propuestos.\\

Este proyecto es software libre, está liberado con la licencia \cite{gplv3} y se puede encontrar en un repositorio público de GitHub\footnote{\url{https://github.com/jlgallego99/TropesToGo}}.


	% Descripción del problema y hasta donde se llega
	\chapter{Descripción del problema}



	% Estado del arte
	% 	1. Crítica al estado del arte
	% 	2. Propuesta
	\chapter{Estado del arte}

El software libre y sus licencias \cite{gplv3} ha permitido llevar a cabo una expansión del
aprendizaje de la informática sin precedentes.

	
	\chapter{Planificación}
En este capítulo se especifica la metodología de desarrollo que se ha utilizado para guiar el completo desarrollo del trabajo, indicando la planificación temporal, historias de usuario e hitos que han servido para saber en todo momento el estado del desarrollo, con el objetivo de garantizar la calidad del proyecto y que se cumplan sus objetivos correctamente.

\section{Metodología utilizada}
Todo el proyecto, desde su concepción, se ha planteado siguiendo una metodología de desarrollo ágil que siga los principios del manifiesto ágil \cite{agilemanifesto}. El principal objetivo al seguir una metodología ágil es el de asegurar un desarrollo flexible, ágil, centrado en las necesidades de los usuarios a los que se desarrolla el software y en el que exista una tolerancia al cambio.

\begin{itemize}
    \item El desarrollo es iterativo e incremental, en el cual se definen una serie de hitos (o entregables) que conforman un Producto Mínimamente Viable (PMV). Este PMV conforma una parte de la funcionalidad total que se quiere tener al final para cumplir los objetivos especificados, y los tests se asegurarán de que sea completamente funcional, garantizando su calidad.
    \item El desarrollo está organizado por una serie de historias de usuario que definen de forma precisa y adecuada el funcionamiento del software y cómo organizar el desarrollo del mismo, siempre desde el punto de si satisfacen al usuario, que es para quien se desarrolla.
    \item Durante todo el desarrollo se acepta que surjan problemas o nuevos requisitos, que se documentarán adecuadamente y se incorporarán a la planificación en mitad del desarrollo para adaptarse a esos cambios lo antes posible y siempre cumplir las necesidades de los usuarios.
    \item Se busca siempre la calidad del software, usando para ello las mejores prácticas del lenguaje tanto en código como en su documentación, teniendo que estar todo debidamente testeado. Se va a desarrollar una herramienta que en última instancia usarán usuarios reales y estos deben obtener el mejor producto posible.
\end{itemize}

Este tipo de desarrollo se adapta bien a un proyecto como este, en el que pueden ir surgiendo necesidades conforme se vaya desarrollando el scraper puesto que es un problema del cual no se sabe desde el principio todas las necesidades que va a plantear debido a la complejidad de la web de TvTropes y a la de las propias herramientas de scraping. Por tanto, esta metodología permite que los problemas que puedan surgir en medio del desarrollo se aborden a tiempo, documentándolos correctamente y sabiendo en todo momento qué se ha hecho, cual es el estado del desarrollo del proyecto, qué hacer y cómo hacerlo.

Con el propósito de cumplir correctamente todos los principios ágiles que se han definido se hacen uso de varias herramientas para el seguimiento del desarrollo y que se describen a continuación.

\section{Seguimiento del desarrollo}
\subsection{GitHub}
Al ser un proyecto de software libre el uso de Git y GitHub como software de control de versiones es imprescindible. Esto permite tener un seguimiento cercano de todos los avances que se hacen tanto en el código como en la documentación, que están ambos alojados en el repositorio público\footnote{\url{https://github.com/jlgallego99/TropesToGo}}. Además de esto también facilita el poder recuperar versiones anteriores del software en caso de algún fallo o cambio sustancial, por lo que es ideal en una metodología ágil en la cual existen constantemente cambios.

Sin embargo lo más importante es que se guía todo el proceso de desarrollo mediante el repositorio de GitHub, teniendo en un mismo lugar todo lo necesario para hacer un correcto seguimiento del estado del trabajo que se tiene en cualquier momento mediante las herramientas que proporciona GitHub. El funcionamiento es el siguiente:
\begin{itemize}
    \item Para llevar el seguimiento de las historias de usuario y tareas asociadas a ellas de la metodología de desarrollo ágil se han creado una serie de issues en el repositorio de GitHub que especifican el trabajo que hay que realizar. Estos issues servirán para saber en todo momento en qué trabajar y siempre están orientados a resolver una historia de usuario, es decir, satisfacer sus necesidades.\\
    Cada vez que, en mitad del desarrollo, se necesite desarrollar algo o solucionar un problema se crea un issue que especifica qué se quiere conseguir y entonces se trabaja en cumplirlo.
    \item Los distintos hitos estarán documentados en la sección de milestones, indicando a modo de resumen cual es el objetivo que se pretende alcanzar al tener ese producto funcional. Este es el principal artefacto para conocer el estado del proyecto, ya que cada hito indica una fecha límite para completarlo y una serie de issues necesarios para completarlo por completo. Conforme se van completando las tareas se puede ver cuánto trabajo queda por hacer.
    \item Cada commit es un avance en el código o la documentación, y en el mensaje se hará una pequeña indicación que sirva para saber qué se ha hecho para avanzar en el issue al que hace referencia. De esta manera podemos saber en todo momento las decisiones que se han tomado para resolver una determinada tarea, teniendo un historial con todos los cambios.
    \item Todo el código y documentación que esté en la rama principal del repositorio se entiende que está probado y por tanto es completamente funcional. \\
    El código se desarrolla en ramas separadas a la principal, que hacen referencia a cada hito que se quiere alcanzar. Cada vez que se complete una tarea de desarrollo se hará un pull request a la rama principal, y si pasa los flujos de CI configurados (los cuales se explicarán en próximas secciones) entonces se dará la tarea por completada y se tendrá una nueva versión funcional en la rama principal. 
\end{itemize}

Como podemos ver, mediante GitHub se tiene en un mismo sitio todo lo necesario para el desarrollo: código, documentación y seguimiento del desarrollo para saber siempre el trabajo que se tiene que hacer y poder avanzar de una forma más cómoda en alcanzar los objetivos y satisfacer al usuario.

\subsection{Hitos}
Los hitos o milestones representan cada uno de los estados en los que va a ir evolucionando la aplicación, conformando un Producto Mínimamente Viable. Cada PMV se construye encima de los anteriores, por lo que no se avanza en el desarrollo hasta que un hito quede completo en su totalidad. Este se considera como completado una vez se hayan resuelto todas sus tareas asociadas, o issues de GitHub.

\subsection{Desarrollo dirigido por pruebas}

\subsection{Integración continua}

\section{Usuarios o partes interesadas}
Antes de definir las historias de usuario, se han analizado los usuarios o partes interesadas en el proyecto, para enfocar el proyecto en resolver sus problemas. Son los siguientes:
\begin{itemize}
    \item \textbf{Tribunal}: el tribunal es el encargado de valorar este trabajo de fin de máster y, por tanto, busca que tenga una memoria bien redactada con todo el proceso de estudio, análisis y desarrollo del proyecto.
    \item \textbf{Investigador}: el investigador, o analista de datos, necesita tener la información estructurada de los tropos para construir modelos que le permitan hacer inferencias sobre las relaciones que tienen los tropos entre sí y llegar a conclusiones que sirvan para su investigación.
\end{itemize}

\section{Historias de usuario}
Sabiendo el problema que se quiere resolver, y las partes interesadas, se pueden definir las historias de usuario que guiarán todo el desarrollo de principio a fin, guiando las decisiones que se irán tomando para llegar a cumplir los objetivos.

\subsection{[HU01] Investigador - Obtener información estructurada sobre todos los tropos}
Como investigador necesito poder obtener información estructurada sobre todos los tropos existentes de cualquier tipo de medio audiovisual, en un formato estandarizado y con una estructura definida para que sean de fácil acceso tanto a humanos como programas para su análisis.

\subsection{[HU02] Investigador - Obtener un subconjunto de la información}
Como investigador necesito poder elegir qué parte de la información sobre tropos quiero obtener, ya sea de un medio audiovisual concreto, todo un género, etc.

\subsection{[HU03] Investigador - Relacionar medios audiovisuales con su ficha en IMDB}
Como investigador necesito obtener información de metadatos sobre cualquier medio audiovisual, para lo cual necesito que los medios audiovisuales de los cuales se han sacado sus tropos se puedan relacionar con su ficha en IMDB.

\subsection{[HU04] Investigador - Información actualizada}
Como investigador quiero que la información de tropos que obtenga esté lo más actualizada posible, sepa cuándo se ha extraído y pueda actualizarla cuando quiera para traer la nueva información que pueda existir.

\section{Temporización}
Una vez vistos los hitos que se esperan alcanzar en el desarrollo, los cuales conformarían cada uno un sprint con una fecha de inicio y de fin, podemos representar la temporización de todo el desarrollo en un diagrama de Gantt.

A lo largo de todo el desarrollo se ha ido manteniendo contacto con el tutor tanto en tutorías presenciales como mediante mensajes por Telegram y GitHub para transmitir y resolver todas las dudas que han ido surgiendo, tanto del código como del propio documento.

	% Análisis del problema
	% 1. Análisis de requisitos
	% 2. Análisis de las soluciones
	% 3. Solucion propuesta
	% 4. Análisis de seguridad
	\chapter{Análisis del problema}
 


	% Desarrollo bajo sprints: 
	% 	1. Permitir registros y login de usuarios
	% 	2. Desarrollo del sistema de incidencias
	% 	3. Desarrollo del sistema de denuncias administrativas y accidentes
	% 	4. Desarrollo del sistema de croquis
	%   5. Instalación de la aplicación de manera automática
	\chapter{Implementación}

La implementación del software se ha dividido en hitos. Estos han sido definidos en GitHub y cada uno de ellos contiene un grupo de \textit{issues} que se corresponden con las distintas mejoras que se han ido incorporando al software a lo largo de su desarrollo.

	% Presupuesto

	% Conclusiones
	\chapter{Conclusiones y trabajos futuros}



	% Trabajos futuros


	
	\newpage
	\bibliography{bibliografia}
	\bibliographystyle{plain}
	
\end{document}

